\documentclass[conference]{IEEEtran}
\IEEEoverridecommandlockouts

\usepackage[T1]{fontenc}
\usepackage{amsmath}
\usepackage{graphicx}
\usepackage{booktabs}
\usepackage{adjustbox}
\usepackage{placeins}
\usepackage{url}
\usepackage[hidelinks]{hyperref}

\graphicspath{{figures/}}

% one plot with its own caption, several of these can share a float
\newcommand{\plot}[3]{%
  \includegraphics[width=\columnwidth]{#1}%
  \caption{#2}\label{#3}}

\begin{document}

\title{CS683 Programming Assignment II:\\
Stride Pattern Prefetching and Instruction-Conscious L2 Replacement}

\author{
\IEEEauthorblockN{Sri Charan Raj Kanakala}
\IEEEauthorblockA{23B1229}
\and
\IEEEauthorblockN{Dheeraj Thammaneni}
\IEEEauthorblockA{23B0997}
\and
\IEEEauthorblockN{Varri Shyam Sri Vardhan}
\IEEEauthorblockA{23B0909}
\and
\IEEEauthorblockN{Srivathsa Thotakura}
\IEEEauthorblockA{23B1305}
}

\maketitle

\begin{abstract}
We implement and evaluate, in ChampSim, (i) a Stride Pattern Prefetcher (SPP)
at the L1D that separates uniform and non-uniform per-IP stride patterns and
throttles its degree by measured accuracy, (ii) SRRIP and an
instruction-conscious I-RRIP policy at the L2C, (iii) a dynamic I-RRIP that
uses leader sets to decide when instruction prioritisation pays off, and
(iv) a prefetch-aware version of dynamic I-RRIP combined with SPP. SPP beats
IP-stride on every trace (geometric-mean speedup 1.0091 vs.\ 1.0023) with a
675-byte budget. I-RRIP cuts L2 instruction MPKI by up to 24\%, and the
combination of SPP and prefetch-aware dynamic I-RRIP gives the highest
geometric-mean speedup of all configurations (1.0152).
\end{abstract}

\section{Methodology}
All results use a single-core ChampSim configuration with a hashed perceptron
branch predictor, a 48\,KB 12-way L1D (64 sets), a 512\,KB 8-way L2C
(1024 sets) and a 2\,MB 16-way LLC. Every run uses 25M warmup and 25M
simulated instructions on the three provided traces. Speedup is
$\mathrm{IPC}/\mathrm{IPC}_{\mathrm{baseline}}$, where the baseline is no
prefetcher (Task~1, Task~3) or LRU at the L2C (Task~2). L1D coverage is
\begin{equation*}
\text{Coverage} = \frac{\text{L1D load misses}_{\text{base}} - \text{L1D load misses}_{\text{pf}}}{\text{L1D load misses}_{\text{base}}}.
\end{equation*}
Accuracy is ChampSim's useful prefetches divided by prefetches issued to the
lower level. The ``GM'' bar in each speedup plot is the geometric mean over
the three traces.

All three traces have a large code footprint: with LRU, L2 instruction MPKI
(10.4, 20.1 and 1.1) is larger than or comparable to L2 data MPKI (7.3, 8.5
and 2.8). This is why data prefetching alone gives modest speedups, and
why instruction-conscious replacement matters, especially on trace2.

%======================================================================
\section{Task 1: Data Prefetching}

\subsection{Task 1A: The IP-stride prefetcher}
The provided IP-stride prefetcher keeps 24 fully-associative trackers with
LRU replacement. Each tracker stores the full IP, the last cache-line address
and the last stride. On an access the stride to the previous line is
computed; if it equals the stored stride (seen twice in a row), the
prefetcher issues $\mathit{addr}+k\cdot\mathit{stride}$ for $k=1..3$
(fixed degree 3), stopping at the 4\,KB page boundary. Lines are filled into
the L1D when fewer than half the MSHRs are busy and into the L2C otherwise.
It only captures a single constant stride per IP, so a repeating sequence
such as $(+1,+2,+2,+2,+3)^n$ never produces two equal consecutive strides
for long and is mostly missed, and its degree is fixed regardless of how
useful its prefetches are.

\subsection{Task 1B: SPP design}
\textbf{IP table.} 64 direct-mapped entries indexed by the low 6 IP bits with
a 10-bit partial tag. Each entry stores the low 16 bits of the last line
address (enough because strides are limited to one page), the last stride,
a 2-bit uniform-stride confidence and a 12-bit stride signature.

\textbf{Dual pattern classification.} The uniform confidence is incremented
when the new stride equals the last stride and decremented otherwise. An IP
is treated as \emph{uniform} only when the counter is saturated at 3; a
pattern like $1,2,2,2,3$ therefore does not look uniform during its short
$2,2,2$ run. Uniform IPs are checked first and prefetch
$\mathit{line}+k\cdot\mathit{stride}$ for $k=1..\mathit{degree}$. Otherwise the
IP falls back to the non-uniform path.

\textbf{Non-uniform path.} The signature is updated as
$\mathit{sig} \leftarrow ((\mathit{sig}\ll 3) \oplus \mathit{stride}) \bmod 2^{12}$,
so it summarises the last four strides. A 256-entry pattern table, indexed by
$(\mathit{sig} \oplus (\mathit{sig}\gg 8)) \bmod 256$, maps a signature to the
next expected stride with a 2-bit confidence. On every access the entry for
the \emph{old} signature is trained with the stride that actually happened
(increment on a match, otherwise decrement and replace the stride once the
confidence reaches 0). To prefetch, SPP looks up the current signature,
adds the predicted stride, folds that stride into a speculative signature and
looks up again, up to \emph{degree} times. It stops when an entry has
confidence below 2 or the next line leaves the page. For the stride sequence
$(+1,+2,+2,+2,+3)^n$, once trained every signature has a single successor,
so SPP walks the pattern several strides ahead.

\textbf{Page boundary and fill level.} Every candidate is checked against the
4\,KB page of the triggering access. Strides larger than 63 lines reset the
entry since they cannot stay within the page. Like IP-stride, prefetches fill
the L1D when the MSHR is less than half full and the L2C otherwise.

\textbf{Storage overhead.} Table~\ref{tab:storage} gives the exact bit-level
breakdown: 5397 bits, about 675 bytes.

\begin{table}[t]
\centering
\caption{SPP storage overhead}
\label{tab:storage}
\footnotesize
\begin{tabular}{lrr}
\toprule
Structure / field & Bits & Total \\
\midrule
\multicolumn{3}{l}{\emph{IP table, 64 entries $\times$ 48 bits}} \\
\quad valid & 1 & \\
\quad IP tag & 10 & \\
\quad last line address (low bits) & 16 & \\
\quad last stride ($-63..63$) & 7 & \\
\quad uniform confidence & 2 & \\
\quad signature & 12 & \\
\quad \emph{subtotal} & 48 & 3072 \\
\midrule
\multicolumn{3}{l}{\emph{Pattern table, 256 entries $\times$ 9 bits}} \\
\quad next stride & 7 & \\
\quad confidence & 2 & \\
\quad \emph{subtotal} & 9 & 2304 \\
\midrule
\multicolumn{3}{l}{\emph{Accuracy throttling}} \\
\quad prefetches sent in epoch (0..256) & 9 & \\
\quad useful prefetches in epoch (0..256) & 9 & \\
\quad current degree (1..4) & 3 & \\
\quad \emph{subtotal} & & 21 \\
\midrule
\textbf{Total} & & \textbf{5397 bits $\approx$ 675 B} \\
\bottomrule
\end{tabular}
\end{table}

\textbf{Accuracy throttling.} SPP counts prefetches sent to the lower level
and how many of them turned out useful. Every epoch of 256 sent prefetches it
computes the epoch accuracy and sets the degree for the next epoch:
\begin{center}
\footnotesize
\begin{tabular}{lcccc}
\toprule
Accuracy & $<25\%$ & $25$--$50\%$ & $50$--$75\%$ & $\geq 75\%$ \\
Degree & 1 & 2 & 3 & 4 \\
\bottomrule
\end{tabular}
\end{center}

\emph{How the thresholds were chosen.} An epoch of 256 prefetches is long
enough that one prefetch changes the accuracy by only 0.4\%, so the decision
is not driven by noise, yet short enough to follow program phases (a few
thousand loads). It also keeps both epoch counters at 9 bits. The accuracy
range is split into equal quarters because, at degree $d$ and accuracy $a$,
roughly $d\cdot a$ prefetches per trigger are useful and $d(1-a)$ only pollute
the L1D and consume bandwidth. Each step up in degree is only allowed once the
accuracy band guarantees that most of the extra prefetches will be useful.

\emph{Why these degrees.} At $\geq 75\%$ the pattern is clearly regular, and
looking 4 lines ahead hides more of the memory latency. With accuracy between
25 and 75\% we keep a moderate degree of 2--3, which still gives timeliness
for the common case. The prefetchers on these traces sit around 45--55\%
accuracy, where the final degree settled at 3, 2 and 4 on the three traces,
showing the degree moving between bands as phases change. Below 25\% most
prefetches are wasted, so we drop to degree 1 instead of 0: keeping a trickle
of prefetches lets the accuracy be measured again and the degree recover in a
later epoch. The maximum of 4 matches the number of strides captured in the
12-bit signature; walking further ahead mostly reuses speculative signatures
and loses accuracy.

\subsection{Task 1 results}
Table~\ref{tab:task1} and Figs.~\ref{fig:t1-speedup}--\ref{fig:t1-cov} show
the results. SPP outperforms IP-stride on every metric and every trace.
Speedups are 1.0069, 1.0052 and 1.0152 (GM 1.0091) against 1.0039, 1.0007 and
1.0022 for IP-stride. Coverage rises from 5.2/0.7/2.5\% to 8.0/2.3/11.9\%,
while accuracy also increases (52.8$\rightarrow$55.7\%,
38.1$\rightarrow$49.5\%, 15.9$\rightarrow$46.8\%). The largest gain is on
trace3, where IP-stride's accuracy is only 15.9\% because most strides are not
constant. SPP's counters show that non-uniform prefetches outnumber uniform
ones on all traces (e.g.\ 594\,K vs.\ 137\,K on trace3), so the pattern table is
what produces the extra coverage. Absolute speedups remain small because a
large part of the stall time comes from L2 instruction misses, which an L1D
prefetcher cannot remove.

\begin{table}[t]
\centering
\caption{L1D prefetchers: performance and prefetch quality (25M instructions)}
\label{tab:task1}
\footnotesize
\setlength{\tabcolsep}{3pt}
\begin{adjustbox}{max width=\linewidth}
\begin{tabular}{llrrrrr}
\toprule
Trace & Config & IPC & Speedup & L1D MPKI & Acc. (\%) & Cov. (\%) \\
\midrule
trace1 & No prefetcher & 0.2321 & 1.0000 & 10.679 & -- & -- \\
 & IP-stride & 0.2330 & 1.0039 & 10.119 & 52.75 & 5.24 \\
 & SPP & 0.2337 & 1.0069 & 9.824 & 55.65 & 8.01 \\
\midrule
trace2 & No prefetcher & 0.2499 & 1.0000 & 13.535 & -- & -- \\
 & IP-stride & 0.2501 & 1.0007 & 13.437 & 38.10 & 0.72 \\
 & SPP & 0.2512 & 1.0052 & 13.219 & 49.52 & 2.34 \\
\midrule
trace3 & No prefetcher & 0.3820 & 1.0000 & 3.403 & -- & -- \\
 & IP-stride & 0.3828 & 1.0022 & 3.319 & 15.88 & 2.47 \\
 & SPP & 0.3878 & 1.0152 & 2.999 & 46.77 & 11.89 \\
\bottomrule
\end{tabular}
\end{adjustbox}
\end{table}


\begin{figure}[!t]
\centering
\plot{task1/speedup}{Task 1 speedup over no prefetching.}{fig:t1-speedup}
\vspace{2mm}
\plot{task1/mpki}{Task 1 L1D load MPKI.}{fig:t1-mpki}
\end{figure}

\begin{figure}[!t]
\centering
\plot{task1/accuracy}{Task 1 L1D prefetch accuracy.}{fig:t1-acc}
\vspace{2mm}
\plot{task1/coverage}{Task 1 L1D miss coverage.}{fig:t1-cov}
\end{figure}

\FloatBarrier
%======================================================================
\section{Task 2A: Instruction-Conscious L2 Replacement (I-RRIP)}

\subsection{Design}
The baseline SRRIP uses a 2-bit re-reference prediction value (RRPV) per line:
fills are inserted at RRPV 2, hits promote to 0, and the victim is the first
line at RRPV 3 (aging the whole set until one exists). I-RRIP adds one bit per
line recording whether it holds instructions (from the packet's
\texttt{instruction} field, passed to \texttt{update\_replacement\_state}) and
changes all three parts of the policy:

\textbf{Insertion.} Instruction lines are inserted at RRPV 0, data lines at
RRPV 2 as in SRRIP. A code line that missed in the L1I and the L2 will be
fetched again soon if the code footprint exceeds the L1I, and an L2 instruction
miss stalls the whole frontend, whereas the out-of-order core can overlap
many data misses. Inserting code at 0 gives it the longest protection before
it can be aged out. Data keeps SRRIP's insertion so that scans and streaming
data still leave quickly.

\textbf{Promotion.} An instruction hit goes straight to RRPV 0. A data hit
only decrements its RRPV by one. Data therefore has to prove reuse over
several hits before it is as protected as code, so a data line touched once or
twice cannot push hot code out.

\textbf{Eviction.} The victim search first looks for a \emph{data} line at
RRPV 3 and evicts code at RRPV 3 only when no such data line exists; if neither
exists, the set is aged as in SRRIP. Code is never made unevictable: aging
still brings it to RRPV 3, so a dead code region is eventually replaced and
the policy cannot fill a set with stale instructions.

Changing all three parts is what makes the priority effective. Changing only
insertion would let a single data hit promote data to 0, equal to code.
Changing only eviction would still insert code at RRPV 2, where it is aged out
quickly between frontend re-visits. The cost is 1 bit per line (8\,K bits for
the 512\,KB L2C) on top of SRRIP's 2-bit RRPV.

\subsection{Results}
Table~\ref{tab:task2a} and Figs.~\ref{fig:t2a-speedup}--\ref{fig:t2a-data} show
the results. I-RRIP reduces L2 instruction MPKI on every trace, from 10.37 to
9.88 (trace1), 20.14 to 15.30 (trace2, $-24\%$) and 1.12 to 1.00 (trace3). As
expected, data MPKI rises because data lines now lose to code (e.g.\ 7.33 to
7.90 on trace1). On trace2, where code is about 70\% of the L2 misses, this
trade is clearly worth it: total L2C MPKI drops from 28.6 to 24.6 and IPC
improves by 1.86\%. On trace1 the instruction and data changes cancel out
(0.9988), and trace3 has very few instruction misses, so it is unchanged. Plain
SRRIP is within $\pm 0.2\%$ of LRU on all traces. The GM speedup is 1.0058 for
I-RRIP vs.\ 1.0002 for SRRIP. The trace1 slowdown motivates the dynamic policy
in Task 2B.

\begin{table}[t]
\centering
\caption{L2C replacement policies, no L1D prefetcher (25M instructions)}
\label{tab:task2a}
\footnotesize
\setlength{\tabcolsep}{3pt}
\begin{adjustbox}{max width=\linewidth}
\begin{tabular}{llrrrrr}
\toprule
Trace & Policy & IPC & Speedup & L2C MPKI & Instr. MPKI & Data MPKI \\
\midrule
trace1 & LRU & 0.2321 & 1.0000 & 17.702 & 10.371 & 7.331 \\
 & SRRIP & 0.2320 & 0.9996 & 17.608 & 10.550 & 7.059 \\
 & I-RRIP & 0.2319 & 0.9988 & 17.773 & 9.876 & 7.897 \\
\midrule
trace2 & LRU & 0.2499 & 1.0000 & 28.644 & 20.141 & 8.504 \\
 & SRRIP & 0.2504 & 1.0020 & 27.272 & 19.325 & 7.947 \\
 & I-RRIP & 0.2546 & 1.0186 & 24.563 & 15.299 & 9.264 \\
\midrule
trace3 & LRU & 0.3820 & 1.0000 & 3.956 & 1.120 & 2.836 \\
 & SRRIP & 0.3816 & 0.9990 & 4.041 & 1.195 & 2.845 \\
 & I-RRIP & 0.3820 & 1.0001 & 3.889 & 0.995 & 2.894 \\
\bottomrule
\end{tabular}
\end{adjustbox}
\end{table}


\begin{figure}[!t]
\centering
\plot{task2a/speedup}{Task 2A speedup over LRU.}{fig:t2a-speedup}
\vspace{2mm}
\plot{task2a/l2c_mpki}{Task 2A L2C load MPKI.}{fig:t2a-l2c}
\end{figure}

\begin{figure}[!t]
\centering
\plot{task2a/l2_instr_mpki}{Task 2A L2C instruction load MPKI.}{fig:t2a-instr}
\vspace{2mm}
\plot{task2a/l2_data_mpki}{Task 2A L2C data load MPKI.}{fig:t2a-data}
\end{figure}

\FloatBarrier
%======================================================================
\section{Task 2B: Dynamic I-RRIP}

\subsection{Design and decision logic}
\textbf{Leader sets.} Every 32nd set (sets $0, 32, 64, \dots$) is a leader,
giving 32 leaders out of 1024 sets (3.1\%). Leaders always run plain SRRIP, so
their statistics describe the workload without being biased by instruction
prioritisation. Only 32 leaders are needed because the policy estimates a
\emph{ratio}, not which of two policies misses less: with a few thousand
accesses per halving interval, 32 sets already give a stable estimate, while
keeping 97\% of the cache free to use the better policy. Evenly spaced leaders
also sample the whole index range, so a hot region of the address space does
not bias the estimate.

\textbf{Decision logic.} In the leader sets we count demand accesses and
misses, both in total and for instruction lines (writebacks and, for Task 3,
prefetches are excluded). Follower sets use I-RRIP when
\begin{equation*}
\frac{\text{instr.\ accesses} + \text{instr.\ misses}}{\text{total accesses} + \text{total misses}} > T
\end{equation*}
and SRRIP otherwise. Counting misses as well as accesses weights instruction
traffic that actually misses more heavily, which is the traffic that stalls
the frontend. All four counters are halved every 4096 leader accesses, so the
ratio is an exponentially decaying average and the policy follows phase
changes instead of the whole-program mix. The overhead is four counters and
the threshold, on top of I-RRIP's per-line bit.

\subsection{Threshold sensitivity}
We swept $T \in \{0.05, 0.15, 0.25, 0.35, 0.45\}$
(Figs.~\ref{fig:t2b-sens}--\ref{fig:t2b-sens-instr}; full data in
Table~\ref{tab:task2b}). The GM speedups are 1.0050, 1.0050, 1.0047, 1.0045 and
1.0044. A higher threshold makes the followers fall back to SRRIP more often.
On trace1 the share of follower fills done in I-RRIP mode falls from 99\% at
$T{=}0.15$ to 0.07\% at $T{=}0.45$, where the policy becomes identical to SRRIP
(instruction MPKI 10.55). On trace2, the most code-heavy trace, the I-RRIP share
drops from 92\% to 74\%, and speedup falls monotonically from 1.0161 to 1.0146.
Trace3 is best at $T \le 0.15$. We use $T{=}0.15$: it is tied for the best GM,
and it is the largest value that keeps I-RRIP active on trace2 and trace3, which
gives some margin against switching too eagerly on data-heavy phases.

\subsection{Comparison}
Figs.~\ref{fig:t2b-speedup}--\ref{fig:t2b-data} compare dynamic I-RRIP
($T{=}0.15$) with LRU, SRRIP and static I-RRIP. It keeps most of I-RRIP's gain on
trace2 (1.0161 vs.\ 1.0186) and matches it elsewhere. It is slightly below
static I-RRIP because the 3\% leader sets always run SRRIP, and because the ratio
measures how much code there is rather than whether prioritising it lowers the
miss count. On these three traces code is a large share of the traffic, so the
switch mostly selects I-RRIP. The dynamic mechanism matters for data-dominated
phases, which it handles by falling back to SRRIP (as the $T{=}0.45$ point
shows on trace1).

\begin{table}[t]
\centering
\caption{Dynamic I-RRIP threshold sweep against LRU, SRRIP and I-RRIP (25M instructions)}
\label{tab:task2b}
\footnotesize
\setlength{\tabcolsep}{3pt}
\begin{adjustbox}{max width=\linewidth}
\begin{tabular}{llrrrrr}
\toprule
Trace & Policy & IPC & Speedup & L2C MPKI & Instr. MPKI & Data MPKI \\
\midrule
trace1 & LRU & 0.2321 & 1.0000 & 17.702 & 10.371 & 7.331 \\
 & SRRIP & 0.2320 & 0.9996 & 17.608 & 10.550 & 7.059 \\
 & I-RRIP & 0.2319 & 0.9988 & 17.773 & 9.876 & 7.897 \\
 & Dyn. I-RRIP (0.05) & 0.2319 & 0.9988 & 17.763 & 9.900 & 7.863 \\
 & Dyn. I-RRIP (0.15) & 0.2319 & 0.9988 & 17.775 & 9.902 & 7.873 \\
 & Dyn. I-RRIP (0.25) & 0.2320 & 0.9993 & 17.748 & 9.884 & 7.865 \\
 & Dyn. I-RRIP (0.35) & 0.2320 & 0.9993 & 17.745 & 9.889 & 7.856 \\
 & Dyn. I-RRIP (0.45) & 0.2320 & 0.9996 & 17.607 & 10.548 & 7.059 \\
\midrule
trace2 & LRU & 0.2499 & 1.0000 & 28.644 & 20.141 & 8.504 \\
 & SRRIP & 0.2504 & 1.0020 & 27.272 & 19.325 & 7.947 \\
 & I-RRIP & 0.2546 & 1.0186 & 24.563 & 15.299 & 9.264 \\
 & Dyn. I-RRIP (0.05) & 0.2539 & 1.0161 & 24.982 & 15.779 & 9.203 \\
 & Dyn. I-RRIP (0.15) & 0.2539 & 1.0161 & 24.979 & 15.777 & 9.202 \\
 & Dyn. I-RRIP (0.25) & 0.2538 & 1.0154 & 25.118 & 15.921 & 9.197 \\
 & Dyn. I-RRIP (0.35) & 0.2536 & 1.0149 & 25.186 & 16.006 & 9.181 \\
 & Dyn. I-RRIP (0.45) & 0.2536 & 1.0146 & 25.275 & 16.128 & 9.147 \\
\midrule
trace3 & LRU & 0.3820 & 1.0000 & 3.956 & 1.120 & 2.836 \\
 & SRRIP & 0.3816 & 0.9990 & 4.041 & 1.195 & 2.845 \\
 & I-RRIP & 0.3820 & 1.0001 & 3.889 & 0.995 & 2.894 \\
 & Dyn. I-RRIP (0.05) & 0.3820 & 1.0000 & 3.894 & 1.003 & 2.891 \\
 & Dyn. I-RRIP (0.15) & 0.3820 & 1.0000 & 3.891 & 1.002 & 2.890 \\
 & Dyn. I-RRIP (0.25) & 0.3818 & 0.9994 & 3.975 & 1.088 & 2.887 \\
 & Dyn. I-RRIP (0.35) & 0.3817 & 0.9993 & 4.031 & 1.143 & 2.887 \\
 & Dyn. I-RRIP (0.45) & 0.3816 & 0.9990 & 4.049 & 1.164 & 2.885 \\
\bottomrule
\end{tabular}
\end{adjustbox}
\end{table}


\begin{figure}[!t]
\centering
\plot{task2b/sensitivity_speedup}{Task 2B threshold sensitivity: speedup over LRU.}{fig:t2b-sens}
\vspace{2mm}
\plot{task2b/sensitivity_l2_instr_mpki}{Task 2B threshold sensitivity: L2C instruction MPKI.}{fig:t2b-sens-instr}
\end{figure}

\begin{figure}[!t]
\centering
\plot{task2b/speedup}{Task 2B speedup over LRU ($T{=}0.15$).}{fig:t2b-speedup}
\vspace{2mm}
\plot{task2b/l2c_mpki}{Task 2B L2C load MPKI.}{fig:t2b-l2c}
\end{figure}

\begin{figure}[!t]
\centering
\plot{task2b/l2_instr_mpki}{Task 2B L2C instruction load MPKI.}{fig:t2b-instr}
\vspace{2mm}
\plot{task2b/l2_data_mpki}{Task 2B L2C data load MPKI.}{fig:t2b-data}
\end{figure}

\FloatBarrier
%======================================================================
\section{Task 3: SPP with Prefetch-Aware Dynamic I-RRIP}

\subsection{Handling prefetch packets}
When SPP fills into the L2C, its prefetches arrive at the replacement policy
with type \texttt{PREFETCH}. Treating them like demand data has two problems:
they distort the leader-set ratio, and the roughly 50\% of them that are never
used occupy L2 space at the same priority as useful data. The policy handles them
as follows:
\begin{itemize}
\item \textbf{Leader counters ignore prefetches.} Only demand accesses decide
whether code is important. Otherwise a burst of data prefetches would lower the
instruction ratio and switch off instruction prioritisation exactly when the
memory system is busiest.
\item \textbf{Prefetch fills are inserted at RRPV 3.} With SPP's measured
accuracy of 47--56\%, about half of the prefetched lines are dead on arrival.
Inserting them at the eviction position means a useless prefetch is the next
victim in its set and never displaces code or reused data.
\item \textbf{Prefetch hits do not promote.} A prefetch that hits in the L2 says
nothing about demand reuse, so it leaves the RRPV unchanged.
\item \textbf{Demand hits promote normally.} When the core actually uses a
prefetched line, it is promoted like any other data line (to 0 in SRRIP mode, one
step in I-RRIP mode), so useful prefetches are kept.
\end{itemize}
This is in the spirit of prefetch-aware insertion: unproven prefetches get the
lowest priority, but a correct prefetch is not penalised once it has been used.

\subsection{Results}
Table~\ref{tab:task3} and Figs.~\ref{fig:t3-speedup}--\ref{fig:t3-data} show
the results. The combination has the highest GM speedup of all
configurations: 1.0152, against 1.0091 for SPP alone and 1.0050 for dynamic
I-RRIP alone. It outperforms both components on trace2 (1.0232 vs.\ 1.0052 and
1.0161) and trace3 (1.0159 vs.\ 1.0152 and 1.0000). On trace1 it is 0.04\%
below standalone SPP (1.0065 vs.\ 1.0069), because dynamic I-RRIP alone is
slightly negative on trace1 (Task 2A). The gains add up because the two
mechanisms remove different misses: SPP removes L1D data misses, and I-RRIP
removes L2 instruction misses. The combination also has the lowest L2C load
MPKI on every trace (16.99, 24.63 and 3.30). L2 data MPKI with SPP stays well
below the no-prefetch dynamic I-RRIP configuration (e.g.\ 2.31 vs.\ 2.89 on
trace3), showing that the low-priority insertion keeps prefetches from
polluting the L2. SPP's accuracy and coverage are essentially unchanged by the
L2 policy (Figs.~\ref{fig:t3-acc}, \ref{fig:t3-cov}), as expected for an L1D
prefetcher.

\begin{table*}[t]
\centering
\caption{SPP combined with prefetch-aware dynamic I-RRIP (25M instructions)}
\label{tab:task3}
\footnotesize
\setlength{\tabcolsep}{3pt}
\begin{adjustbox}{max width=\linewidth}
\begin{tabular}{llrrrrrrrr}
\toprule
Trace & Config & IPC & Speedup & Acc. (\%) & Cov. (\%) & L1D & L2C & L2 instr. & L2 data \\
\midrule
trace1 & No PF + LRU & 0.2321 & 1.0000 & -- & -- & 10.679 & 17.702 & 10.371 & 7.331 \\
 & SPP + LRU & 0.2337 & 1.0069 & 55.65 & 8.01 & 9.824 & 17.024 & 10.411 & 6.613 \\
 & No PF + Dyn. I-RRIP & 0.2319 & 0.9988 & -- & -- & 10.680 & 17.775 & 9.902 & 7.873 \\
 & SPP + Dyn. I-RRIP & 0.2336 & 1.0065 & 55.60 & 7.95 & 9.830 & 16.986 & 9.890 & 7.096 \\
\midrule
trace2 & No PF + LRU & 0.2499 & 1.0000 & -- & -- & 13.535 & 28.644 & 20.141 & 8.504 \\
 & SPP + LRU & 0.2512 & 1.0052 & 49.52 & 2.34 & 13.219 & 28.680 & 20.449 & 8.231 \\
 & No PF + Dyn. I-RRIP & 0.2539 & 1.0161 & -- & -- & 13.538 & 24.979 & 15.777 & 9.202 \\
 & SPP + Dyn. I-RRIP & 0.2557 & 1.0232 & 49.96 & 2.32 & 13.221 & 24.628 & 15.684 & 8.944 \\
\midrule
trace3 & No PF + LRU & 0.3820 & 1.0000 & -- & -- & 3.403 & 3.956 & 1.120 & 2.836 \\
 & SPP + LRU & 0.3878 & 1.0152 & 46.77 & 11.89 & 2.999 & 3.402 & 1.130 & 2.271 \\
 & No PF + Dyn. I-RRIP & 0.3820 & 1.0000 & -- & -- & 3.403 & 3.891 & 1.002 & 2.890 \\
 & SPP + Dyn. I-RRIP & 0.3881 & 1.0159 & 47.03 & 12.00 & 2.995 & 3.304 & 0.995 & 2.309 \\
\bottomrule
\end{tabular}
\end{adjustbox}
\end{table*}


\begin{figure}[!t]
\centering
\plot{task3/speedup}{Task 3 speedup over no prefetching with LRU.}{fig:t3-speedup}
\vspace{2mm}
\plot{task3/l1d_mpki}{Task 3 L1D load MPKI.}{fig:t3-l1d}
\end{figure}

\begin{figure}[!t]
\centering
\plot{task3/accuracy}{Task 3 L1D prefetch accuracy.}{fig:t3-acc}
\vspace{2mm}
\plot{task3/coverage}{Task 3 L1D miss coverage.}{fig:t3-cov}
\end{figure}

\begin{figure}[!t]
\centering
\plot{task3/l2c_mpki}{Task 3 L2C load MPKI.}{fig:t3-l2c}
\vspace{2mm}
\plot{task3/l2_instr_mpki}{Task 3 L2C instruction load MPKI.}{fig:t3-instr}
\end{figure}

\begin{figure}[!t]
\centering
\plot{task3/l2_data_mpki}{Task 3 L2C data load MPKI.}{fig:t3-data}
\end{figure}

\FloatBarrier
%======================================================================
\section{Conclusion}
SPP's dual uniform / signature-based classification and accuracy throttling
give consistently higher accuracy, coverage and speedup than IP-stride for
675\,B of state. At the L2C, prioritising instructions in insertion, promotion
and eviction cuts instruction MPKI by up to 24\%. The leader-set based dynamic
policy keeps that benefit while falling back to SRRIP when code is a small
share of the traffic. Making the policy prefetch-aware lets the two mechanisms
work together, giving the best overall speedup (GM 1.0152) and the lowest L2C
MPKI on all traces.

\end{document}
